% ============================================================
% Cross-Sector Layer 4 Closure of the Substrate Chirality Handle:
%   Electroweak V--A Coupling and Quark Chiral-Polarity-Bias from |M| = chi/6
% A Joint Conditional Theorem Closure Paper for OPEN-FP-SF-2-CHIR and
%   SM-2 v2.0+ Chiral-Polarity-Bias EFT Continuum-Limit
% Conscious Point Physics Flagship Paper Series (Chirality Continuum line)
% Version history: see flagship_papers/chirality_continuum/README.md
% ============================================================

\documentclass[11pt,letterpaper]{article}
\usepackage[utf8]{inputenc}
\usepackage[margin=1in]{geometry}
\usepackage[T1]{fontenc}
\usepackage{amsmath,amssymb,amsfonts,amsthm}
\usepackage{graphicx}
\usepackage{tikz}
\usetikzlibrary{positioning,calc,arrows.meta,shapes.geometric}
\usepackage{xcolor}
\usepackage{hyperref}
\usepackage{booktabs}
\usepackage{array}
\usepackage{enumitem}
\usepackage{float}
\usepackage[font=small, labelfont=bf]{caption}
\usepackage{parskip}
\usepackage{mdframed}

\hypersetup{
    colorlinks=true,
    linkcolor=blue,
    citecolor=blue,
    urlcolor=blue
}

\newtheorem{theorem}{Theorem}[section]
\newtheorem{proposition}[theorem]{Proposition}
\newtheorem{lemma}[theorem]{Lemma}
\newtheorem{corollary}[theorem]{Corollary}
\newtheorem{conjecture}[theorem]{Conjecture}
\newtheorem{definition}[theorem]{Definition}
\newtheorem{remark}[theorem]{Remark}
\newtheorem{openproblem}{Open Problem}

% Shared notation conventions inherited from Capotauro v2.0
\newcommand{\phig}{\varphi}
\newcommand{\Mzero}{M_0}
\newcommand{\Kthree}{\ensuremath{K_3}}
\newcommand{\DeltapLR}{\Delta p_{LR}}
\newcommand{\Cchi}{\hat{C}_\chi}
\newcommand{\PsiminusOne}{|\Psi_-^{(1)}\rangle}
\newcommand{\PsiminusTwo}{|\Psi_-^{(2)}\rangle}
\newcommand{\chiplus}{|\chi_+\rangle}
\newcommand{\chiminus}{|\chi_-\rangle}
\newcommand{\ZBW}{\mathrm{ZBW}}
\newcommand{\Dsix}{D_6}
\newcommand{\Sthree}{S_3}
\newcommand{\Hthree}{H_3}
\newcommand{\Hfour}{H_4}
\newcommand{\Ih}{I_h}

% Continuum-EFT specific notation
\newcommand{\Phicont}{\Phi_*}
\newcommand{\Oeff}{\mathcal{O}^{\text{eff}}}
\newcommand{\OeffW}{\mathcal{O}^{\text{eff,W}}}
\newcommand{\OeffqDP}{\mathcal{O}^{\text{eff,qDP}}}
\newcommand{\Lambdasub}{\Lambda_{\text{sub}}}
\newcommand{\muobs}{\mu_{\text{obs}}}
\newcommand{\elledge}{\ell_{\text{edge}}}
\newcommand{\DeltaFqDP}{\Delta F^{qDP}}
\newcommand{\FeDPref}{F^{eDP}_{\text{ref}}}
\newcommand{\Vcage}{V_{\text{cage}}}
\newcommand{\dGamma}{d_\Gamma}

\title{Cross-Sector Layer 4 Closure of the Substrate Chirality Handle:\\
Electroweak V--A Coupling and Quark Chiral-Polarity-Bias from $|M| = \chi/6$}

\author{Thomas Lee Abshier, ND\thanks{Hyperphysics Institute / Renaissance Ministries. Email: drthomas@hyperphysics-institute.org} \and
Claude Opus\thanks{Anthropic. Methodological collaborator and co-author. The first author conceived the framework and directed the closure trajectory; the second author contributed mathematical formalization, computation, and documentation under the first author's direction.}}

\date{Version 0.5 (DRAFT) --- 20 May 2026\\
\small{Conscious Point Physics Flagship Paper Series (Chirality Continuum line)}}

\begin{document}
\maketitle

% ============================================================
% ABSTRACT
% ============================================================
\begin{abstract}
We present a joint Layer 4 continuum-EFT closure of two Standard Model phenomenological features --- the V--A coupling structure of $W^\pm$-mediated weak interactions and the chiral-polarity-bias structure of quark charge asymmetry --- as parallel projections of a single substrate-level chirality matrix element $|M| = \chi/6 = \varphi^{-3}/6 \approx 0.0394$. The substrate handle is inherited at full Layer 3 rigor from Capotauro v2.0's cross-sector unification theorems (THEO-SD-CHIR-1 for the W-bracelet sector under the electroweak V--A manifestation; THEO-SD-CHIR-2 for the qDP/eDP sector under the electromagnetic-handedness manifestation), together with the K3-doublet sector closure (THEO-CAP-1) under the OPEN-SD-CHIR-PRIMITIVE umbrella's three-way cross-sector unification at substrate level. The shared substrate-handle-to-effective-coupling bridge step is identified as the load-bearing technical content and closed at sector-agnostic level as Theorem 3.1 (programme-level registration THEO-CHIR-CONT-1); the two sector-specific Layer 4 closures (Theorems 4.4 = THEO-CHIR-CONT-2 and 5.5 = THEO-CHIR-CONT-3) proceed via standard SM machinery (Yang-Mills $SU(2)_L \times U(1)_Y$ effective field theory for the V--A coupling kinematic projection; effective free-energy partition-function framework for the chiral-polarity-bias thermodynamic projection). The two sector-specific projections converge at observable level on a single primary empirical channel --- the leptogenesis CP-asymmetry $\DeltapLR \approx 0.0394$ inferred from the observed baryon asymmetry of the universe at $\eta_B \sim 6 \times 10^{-10}$ via standard thermal-leptogenesis machinery --- providing cross-sector convergent empirical anchor at $\chi/6 \approx 0.0394$ versus observed $\sim 0.04$ (match within 2\% of substrate-handle prediction at current $\sigma \sim 5 \times 10^{-3}$ BAU back-derivation precision). The paper constitutes the second cross-sector closure in Conscious Point Physics after SF-4 v4.0 (Composite K3-Cage-Shell Coupling Theorem) and establishes the THEO-CHIR-CONT-N sub-prefix convention for Layer 4 continuum-EFT projection closures under the OPEN-SD-CHIR-PRIMITIVE umbrella. The closure is conditional theorem closure within the current CPP theorem stack: it rests on the foundational input stack inherited from Capotauro v2.0, on a small number of additional sector-specific continuum-EFT framework foundational inputs, and on the CPP axiom set; first-principles derivation of the substrate primitive 4D direction $\hat{n}$ and of the substrate chirality magnitude $|\chi| = \varphi^{-3}$ remains as future-window work tied to the Q1$'$+Q1$'$.A Layer 3 promotion programme.
\end{abstract}

% ============================================================
% PLAIN-LANGUAGE SUMMARY
% ============================================================
\section*{Plain-language summary}

One number in the Conscious Point Physics framework --- the substrate-level chirality magnitude $\chi/6 = \varphi^{-3}/6 \approx 0.0394$ --- controls two phenomena that ordinary Standard Model physics treats as completely separate. The first is the fact that the weak nuclear force couples to left-handed particles only and not to their mirror-image right-handed counterparts (the so-called V--A structure of weak interactions, established by Goldhaber, Grodzins, and Sunyar in 1958 and refined to sub-percent precision in subsequent decades). The second is the empirical absence of fractional-charge configurations of the kind that would correspond to ``positive down-type quarks'' --- quarks with the wrong sign of electric charge. These two phenomena look unrelated in the textbook Standard Model: V--A is about angular momentum and chirality eigenstates of fermion fields; the chiral-polarity-bias is about which configurations of quark-like substructures actually exist in nature. We show that both are controlled by the same number $\chi/6$, which is itself controlled by a single substrate primitive: a 4D direction $\hat{n}$ in the ambient four-dimensional Euclidean space in which the 600-cell substrate of Conscious Point Physics lives, together with a polytope-geometric magnitude $\varphi^{-3}$ derived from the golden ratio $\varphi = (1+\sqrt{5})/2$ via the perturbative-distance-ratio constraint at the substrate's first-shell icosahedron.

The mechanism is the following. The 600-cell substrate has, at any selected host vertex $v_{\text{host}}$, a residual icosahedral symmetry group $\Ih$ which is partially broken by a substrate primitive 4D direction $\hat{n}$ aligned with $v_{\text{host}}$ itself. The 12 vertices of the host's first-shell icosahedron host three structurally distinct substrate objects --- a K3-doublet triangular face (the seat of mass-mixing chirality), a W-bracelet Petrie hexagon (the seat of electroweak V--A), and Linear-ZBW configurations on $\pm$qCP centers (the seat of electromagnetic handedness) --- and each substrate object carries a sector-specific stabilizer subgroup, a sector-specific pairing-convention generator, and a sector-specific chirality operator. Despite the structural distinctness of the three sector-specific data, the chirality matrix element on each sector is the same number $\chi/6 \approx 0.0394$, because the three sectors all sit on the same 12-vertex icosahedral cage with the same effective matter-doublet dimension 2. This three-way cross-sector identity at the substrate level was established by Capotauro v2.0 (Abshier 2026). The present paper shows that the substrate-level identity is preserved unchanged under the projection from the substrate cutoff $\Lambdasub = \elledge^{-1}$ down to the observable scales of the Standard Model --- that is, the substrate magnitude $\chi/6$ feeds the V--A coupling and the chiral-polarity-bias at Standard Model scales without any renormalization correction or scale-dependent suppression. The projection from substrate to observable scales is the standard machinery of Wilson--Fisher renormalization-group flow applied to the substrate's discrete polytope-geometric structure, and it preserves topological substrate quantities (in a sense made precise in §3) at leading order.

The empirical anchor for the joint closure is the leptogenesis CP-asymmetry $\DeltapLR^{\text{obs}} \sim 0.04$ inferred from the observed matter-antimatter asymmetry of the universe via standard thermal-leptogenesis arguments (Davidson, Nardi, and Nir 2008). Our predicted value $\chi/6 \approx 0.0394$ matches this empirical anchor within $2\%$, which is well within the current $\sim 5 \times 10^{-3}$ uncertainty in the back-derivation from cosmological observations. This is a single empirical observable that simultaneously tests both Layer 4 closures of the joint paper: the V--A coupling closure (Section 4) inherits the substrate-handle magnitude via Yang-Mills effective field theory kinematic projection at the massless helicity limit; the chiral-polarity-bias closure (Section 5) inherits the substrate-handle magnitude via effective free-energy partition-function thermodynamic projection at thermal equilibrium scales. The convergence of both closures on the same empirical observable is a structural prediction of the joint paper format and constitutes the structural payoff of the cross-sector closure approach.

% ============================================================
% §1 INTRODUCTION
% ============================================================
\section{Introduction}
\label{sec:introduction}

\subsection{The two Standard Model phenomenological features this paper closes}
\label{sec:two_features}

This paper provides Layer 4 continuum-effective-field-theory closure of two Standard Model phenomenological features that have, until now, been treated independently in the Conscious Point Physics (CPP) framework: the V--A coupling structure of $W^\pm$-mediated weak interactions, and the chiral-polarity-bias structure of quark charge asymmetry.

The first feature --- the V--A coupling structure of charged-current weak interactions --- is the assertion that the $W^\pm$ bosons couple only to the left-handed (LH) chirality projections of fermion fields, with the right-handed (RH) projections decoupled from $W^\pm$ at observable scales \cite{CommonsBucksbaum, ChengLi, Peskin}. This is the maximally parity-violating coupling structure first established by the experiments of Wu et al. \cite{Wu1957} on $^{60}$Co beta decay and Goldhaber, Grodzins, and Sunyar \cite{Goldhaber1958} on $^{152m}$Eu electron capture, and subsequently refined to sub-percent precision through Michel parameter measurements at PSI and TRIUMF \cite{TWISTrho, Michel1950, MarcianoSirlin1988}, $\tau$-polarization measurements at LEP and SLC \cite{LEPElectroweak}, and top-quark spin-correlation measurements at the LHC \cite{ATLASTopSpin, CMSTopSpin}. The V--A structure is conventionally encoded in the Standard Model Lagrangian through the projection operators $P_L = (1-\gamma_5)/2$ and $P_R = (1+\gamma_5)/2$, with the charged-current Lagrangian taking the form $\mathcal{L}_{\text{CC}} = -(g/\sqrt{2}) W^+_\mu \bar{\psi}_L \gamma^\mu \psi_L + \text{h.c.}$ This structure is parametrically pinned in the general ten-coupling parametrization $\{g^\gamma_{\epsilon\mu}\}_{\gamma \in \{S,V,T\}; \epsilon,\mu \in \{L,R\}}$ at $g^V_{LL} = 1$ with all other couplings vanishing, and is consistent with observation to within $\sim 10^{-3}$ precision at all currently accessible energy scales. The CPP question is: \emph{why} is the coupling pinned at $g^V_{LL} = 1$ rather than at some other value or some other coupling-tensor structure?

The second feature --- the chiral-polarity-bias of quark charge asymmetry --- is the empirical absence of fractional-charge configurations corresponding to ``positive down-type quarks'' (quarks with electric charge $+\frac{1}{3}$ rather than the observed $-\frac{1}{3}$). In the CPP framework's substrate-level description, quarks are modeled as Linear-ZBW (zero-base-vector) configurations attached to $\pm$qCP centers at substrate host vertices \cite{SM2v1}, and the asymmetric stabilization of Linear-ZBW configurations on $-$qCP centers versus $+$qCP centers determines which quark charge signs actually occur in nature. SM-2 v1.0 §10 \cite{SM2v1} formulates this as the chiral-polarity-bias mechanism: ``the 600-cell's intrinsic chirality (activated during the Capotauro symmetry-breaking event) preferentially stabilises linear ZBW extras on negative ($-$qCP) centres.'' At observable scales, this mechanism contributes to a polarization-asymmetry observable $\DeltapLR$ which is identified, via the standard thermal-leptogenesis machinery, with the leptogenesis CP-asymmetry parameter $\epsilon_{CP}$ that sources the baryon asymmetry of the universe (BAU) at $\eta_B = n_B/n_\gamma \sim 6 \times 10^{-10}$ \cite{DavidsonNardiNir}. The CPP question is: \emph{why} is the chirality bias of magnitude $\sim 0.04$ rather than of order unity (full chirality selection) or of order $10^{-10}$ (negligible)?

\subsection{Why these two features are paired: shared substrate handle from Capotauro v2.0}
\label{sec:why_paired}

The reason to address these two features jointly --- rather than in two separate single-sector papers --- is that both inherit the same substrate-level chirality matrix element from the OPEN-SD-CHIR-PRIMITIVE umbrella's three-way cross-sector unification at substrate level, established by Capotauro v2.0 \cite{Capotauro}. Three substrate objects on the host vertex's first-shell icosahedron --- the K3-doublet triangular face (mass-mixing sector), the W-bracelet Petrie hexagon (electroweak V--A sector), and the Linear-ZBW configurations on $\pm$qCP centers (electromagnetic-handedness sector) --- all produce identical chirality matrix element magnitudes at full Layer 3 rigor:
\begin{equation}
\left| M^{K_3} \right| = \left| M^W \right| = \left| M^{qDP} \right| = \frac{\chi}{6} = \frac{\varphi^{-3}}{6} \approx 0.0394.
\label{eq:three_way_unification}
\end{equation}
This three-way identity is established by three programme-level theorems registered in the CPP theorem stack: THEO-CAP-1 (theorem \#62; Composite Capotauro Wigner-Eckart Theorem; K3-doublet sector \cite{Capotauro, AbshierCPP2026}); THEO-SD-CHIR-1 (theorem \#63; Cross-Sector Substrate Chirality Unification Theorem; W-bracelet sector \cite{Capotauro}); THEO-SD-CHIR-2 (theorem \#64; qDP/eDP Sector Substrate Chirality Closure Theorem; qDP/eDP sector \cite{Capotauro}). The single substrate parameter $\chi = \varphi^{-3}$ feeds three structurally distinct sectors via three cage-shell averaging operations on the same 12-vertex icosahedral cage with the same effective matter-doublet dimension $\dGamma = 2$.

What Capotauro v2.0 \emph{did not} close is the projection from the substrate-level matrix element to observable phenomena at Standard Model accessible scales. The substrate-level identity \eqref{eq:three_way_unification} lives at the substrate cutoff $\Lambdasub = \elledge^{-1}$, where $\elledge$ is the 600-cell edge length at the substrate's Planck-scale geometric resolution. The two phenomenological features described in §\ref{sec:two_features} live at observable kinematic scales $\muobs \ll \Lambdasub$: the V--A coupling at the electroweak scale $\muobs^W \sim 100$ GeV; the chiral-polarity-bias at thermal-equilibrium scales spanning the QCD scale ($\sim 100$ MeV) through the leptogenesis era ($\sim 10^{12}$ GeV). Bridging the substrate handle at $\Lambdasub$ to the observable kinematic scale at $\muobs^{\text{sector}}$ is the load-bearing Layer 4 technical content that this paper closes for both sectors jointly.

The structural efficiency of the joint paper format derives from the fact that the substrate-handle-to-effective-coupling bridge step is sector-agnostic by construction: the bridge theorem (§\ref{sec:bridge}, Theorem 3.1; programme-level registration THEO-CHIR-CONT-1) depends only on universal substrate-level data $(|\chi|, \dGamma/\Vcage) = (\varphi^{-3}, 1/6)$ and not on sector-specific stabilizers, $\zeta$ generators, or chirality operators. Doing the bridge work once at sector-agnostic level, then applying it twice at sector-specific level, saves an estimated 4--11 sessions of substantive derivation effort versus the fallback of two separate single-sector papers (see §\ref{sec:cross_sector_unification} for the empirical validation of this projection). More importantly, the joint paper format makes the cross-sector convergence at observable level --- both sectors converging on the same empirical observable, the leptogenesis CP-asymmetry $\DeltapLR \approx 0.0394$ --- a structural prediction of the same bridge theorem rather than an emergent empirical coincidence under separate-paper closure.

\subsection{What this paper delivers}
\label{sec:what_delivered}

The paper delivers theorem-level closure of two open foundational-physics entries in the CPP programme's research-frontier ledger \cite{ResearchFrontier}:

\begin{itemize}[leftmargin=*]
\item \textbf{OPEN-FP-SF-2-CHIR}: Layer 4 closure of the SF-2 v2.0+ V--A coupling derivation at the massless helicity limit, including pinning of the coupling to $g^V_{LL} = 1$, derivation of the Michel parameter $\rho = 3/4$ at finite mass via standard V--A four-fermion kinematics, derivation of the 100\% LH preference at the massless helicity limit via chirality-helicity coincidence, and activation of Capotauro Falsifier 6 at three observable thresholds. Closure is at theorem-level rigor as Theorem 4.4 (programme-level registration THEO-CHIR-CONT-2; \cite{ResearchFrontier} theorem \#66).

\item \textbf{SM-2 v2.0+ chiral-polarity-bias EFT continuum-limit closure}: Layer 4 closure of the SM-2 v2.0+ chiral-polarity-bias mechanism at observable thermodynamic scales, including identification of the continuum-EFT operator as the chirality-asymmetric stabilization-energy operator $\DeltaFqDP$ in effective free-energy partition-function framework, derivation of the substrate-level stabilization energy magnitude $|M^{\text{eff,qDP}}| = \chi/6$ at leading order via topological-projection argument, derivation of the exclusion bound at observable thermodynamic scales $\DeltapLR \approx \chi/6 \approx 0.0394$ via Boltzmann-like thermodynamic distribution, and cross-validation against SM-2 v1.0 §10 substrate-level mechanism and against the Sector A V--A coupling derivation. Closure is at theorem-level rigor as Theorem 5.5 (programme-level registration THEO-CHIR-CONT-3; \cite{ResearchFrontier} theorem \#67).
\end{itemize}

The two sector-specific closures share a common bridge step at sector-agnostic level (the substrate-handle-to-effective-coupling bridge; §\ref{sec:bridge}, Theorem 3.1; programme-level registration THEO-CHIR-CONT-1; \cite{ResearchFrontier} theorem \#65). The cross-sector unification at observable level (§\ref{sec:cross_sector_unification}) is the structural identity claim of the paper: a single substrate primitive (the 4D direction $\hat{n}$) with derived magnitude $|\chi| = \varphi^{-3}$ controls every parity-sensitive observable in the CPP framework via shared three-step machinery (substrate-locality + cage-shell averaging + sector-specific pairing convention) at all four levels (Layer 3 substrate-level, Layer 4 sector-agnostic continuum-EFT, Layer 4 sector-specific continuum-EFT, observable scale).

\subsection{Closure status and roadmap}
\label{sec:closure_status}

This paper is a \emph{conditional theorem closure paper} within the current CPP theorem stack \cite{ConditionalClosure}. The closures presented in §\ref{sec:bridge}, §\ref{sec:sector_a}, and §\ref{sec:sector_b} are conditional on a foundational input stack inherited from Capotauro v2.0 \cite{Capotauro} plus a small number of additional sector-specific continuum-EFT framework foundational inputs introduced in §\ref{sec:bridge}, §\ref{sec:sector_a}, and §\ref{sec:sector_b}; on the CPP axiom set \cite{AxiomDocument}; and on the proof outlines of the Yang-Mills continuum-EFT framework anchored at SF-2 v1.0 \cite{SF2v1, EW5} and the effective free-energy partition-function framework anchored at SM-2 v1.0 \cite{SM2v1}.

The paper inherits Capotauro v2.0's conditional-closure discipline. Open verification items deferred to future work are explicitly catalogued in §\ref{sec:open_work}: (i) first-principles derivation of the substrate primitive 4D direction $\hat{n}$ from CPP primitive axioms (the Q1$'$+Q1$'$.A Layer 3 promotion programme; the dynamical-substrate-law gate identified by external reviewers as the defining next gate for the Capotauro programme \cite{ReviewerNotes}); (ii) first-principles derivation of the substrate chirality magnitude $|\chi| = \varphi^{-3}$ from CPP primitive axioms (currently Layer 2 via the perturbative-distance-ratio constraint plus Capotauro v2.0 §sec:chi\_resolution); (iii) Picture A alternative continuum-EFT framework parametrization (OPEN-FP-SF-4-1 candidate; complementary closure to the Picture B Wigner-Eckart EFT framework completed in this paper); (iv) Layer 4 promotion of OPEN-SD-CHIR-PRIMITIVE umbrella manifestations (iv) thermodynamic causal arrow and (v) cosmological-vacuum asymmetry (templates established by THEO-CHIR-CONT-N convention for future work via THEO-CHIR-CONT-4 / -5 candidates); (v) future-collider precision improvements on the three Capotauro Falsifier 6 thresholds (Michel, massless-helicity, leptogenesis CP-asymmetry) at $\sim 10^{-3}$ to $\sim 10^{-4}$ level via CMB-S4 + LiteBIRD + LEGEND-1000 + nEXO + CUPID (leptogenesis); MEG-II + FCC-ee (Michel); CLIC + ILC + LHC high-luminosity (massless-helicity); FCC-ee Higgs program (combined).

The roadmap of the paper is as follows. §\ref{sec:inheritance} establishes the substrate-level inheritance from Capotauro v2.0: the three-way cross-sector unification at substrate level under the OPEN-SD-CHIR-PRIMITIVE umbrella, the conditional-closure foundational input stack, and the Capotauro v2.0 axiom set load-bearing for the joint paper. §\ref{sec:bridge} closes the substrate-handle-to-effective-coupling bridge at sector-agnostic level via Theorem 3.1 (THEO-CHIR-CONT-1), with three sub-statements as named lemmas / sub-theorems (Lemma 3.1 / THEO-CHIR-CONT-1.1; Theorem 3.2 / THEO-CHIR-CONT-1.2; Theorem 3.3 / THEO-CHIR-CONT-1.3). §\ref{sec:sector_a} closes Sector A (electroweak V--A coupling derivation) via Theorem 4.4 (THEO-CHIR-CONT-2) with four sub-claim consequences (sector-specific operator identification; Michel parameter at finite mass; massless-helicity-limit 100\% LH preference; Capotauro Falsifier 6 activation). §\ref{sec:sector_b} closes Sector B (chiral-polarity-bias EFT continuum-limit) via Theorem 5.5 (THEO-CHIR-CONT-3) with four sub-claim consequences (sector-specific operator identification; substrate-level stabilization energy; exclusion bound at observable thermodynamic scales; SM cross-validation). §\ref{sec:cross_sector_unification} synthesizes the cross-sector unification as the structural identity claim of the paper and validates the joint paper format's structural efficiency. §\ref{sec:predictions} catalogues the paper's zero-parameter predictions and falsifiers (six falsifiers per the v0.1 outline). §\ref{sec:open_work} catalogues open theorem-level work deferred to future patches. §\ref{sec:discussion} discusses the paper's contribution to the CPP programme's cross-sector closure methodology and its implications for the OPEN-SD-CHIR-PRIMITIVE umbrella's remaining two observable manifestations.

% ============================================================
% §2 INHERITANCE FROM CAPOTAURO v2.0
% ============================================================
\section{Inheritance from Capotauro v2.0: substrate-level results}
\label{sec:inheritance}

This section establishes the substrate-level results inherited from Capotauro v2.0 \cite{Capotauro} that serve as the substrate-handle foundation for the joint paper's Layer 4 closure work. The inheritance is at theorem-level rigor: Capotauro v2.0 v1.0 SHIPPED status (\cite{Capotauro}; Session 135 Patch 0479) confirms that the substrate-level theorems registered in the CPP theorem stack \cite{ResearchFrontier} have paper-level publication venue and are available for inheritance citation in subsequent Layer 4 work.

\subsection{THEO-SD-CHIR-1 inheritance: K3-doublet $\leftrightarrow$ W-bracelet substrate-level cross-sector unification}
\label{sec:theo_sd_chir_1}

THEO-SD-CHIR-1 \cite{ResearchFrontier, Capotauro} establishes at full Layer 3 rigor that the composite chirality matrix elements on the K3-doublet (mass-mixing sector) and the W-bracelet (electroweak V--A sector) are identical:
\begin{equation}
\left| M^{K_3} \right| = \left| M^W \right| = \frac{\chi}{6} = \frac{\varphi^{-3}}{6} \approx 0.0394
\label{eq:k3_w_unification}
\end{equation}
under vertex-aligned Reading C ($\hat{n} = v_{\text{host}}$ a 600-cell vertex direction) with perturbative-distance-ratio constraint $\epsilon = \varphi^{-3}$ at substrate level. The four-step proof chain establishes:

\begin{enumerate}[leftmargin=*]
\item \textbf{Local-$\Ih$-preservation theorem} (Capotauro v2.0 \cite{Capotauro} §sec:substrate\_locality; Finding C-W39): the first-shell icosahedron at $v_{\text{host}}$ is preserved as $\Ih$-symmetric at first order in $\epsilon$ with only isotropic radial scaling by $(1 - \epsilon/(2\varphi))$; all first-shell-to-first-shell 600-cell edges have $\hat{e}_{ab} \cdot \hat{n} = 0$ identically in 4D (numerically verified to machine precision: all 30 such edges have $|\hat{e} \cdot \hat{n}| < 10^{-9}$). Substrate chirality magnitude identified directly: $\chi \equiv \epsilon = \varphi^{-3}$ at substrate level.

\item \textbf{Substrate-Locality Unification} (Capotauro v2.0 \cite{Capotauro} §sec:substrate\_locality\_unification; Finding C-W40): the local-$\Ih$-preservation theorem applies uniformly to any first-shell-vertex substrate object on the host vertex's first-shell icosahedron. Both the K3-base (3-vertex triangular face per Finding C-W37) and the W-bracelet (6-vertex Petrie hexagon per Finding C-W36) inherit substrate chirality from the same identification $\chi \equiv \epsilon = \varphi^{-3}$ via cage-shell averaging on respective $\Dsix$ sub-stabilizers of the same $\Hthree = \Ih$ residual symmetry.

\item \textbf{Cage-shell factor identity} (Capotauro v2.0 \cite{Capotauro} §sec:cage\_shell): Schur orthogonality cage-shell averaging on the icosahedral cage shared with K3 gives $|M_\perp^W| = \dGamma/\Vcage = 2/12 = 1/6$, identical to the K3-doublet's cage-shell factor of THEO-CAP-1. The integer values $\dGamma = 2$ (matter-doublet dimension) and $\Vcage = 12$ (icosahedral vertex count) are representation-theoretic and polytope-topological invariants respectively.

\item \textbf{Pairing-convention identification} (Capotauro v2.0 \cite{Capotauro} §sec:pairing\_convention; Finding C-W43): the W-bracelet's $\mathbb{Z}_2$ generator $\zeta^W = r^3$ of dihedral $D_6$ is identified as the icosahedral-center inversion in 4D ambient ($p \mapsto \varphi\hat{n} - p$); linear part $-I$ flips $\hat{n}$ (chirality-flipping) and maps each first-shell vertex to its first-shell antipode. The matter-doublet basis on 3 antipodal pairs spans the 2D $E$-irrep of $\Sthree$, parallel to K3's matter-doublet structure. Chirality operator $\Cchi^W \in B_2(D_6)$ sourced from $\hat{n}$ has non-vanishing matrix element via Wigner-Eckart, factorizing as amplitude factor $|M_{\text{amp}}^W| = \chi$ (chirality-eigenvalue matching analog of K3) times cage-shell factor $|M_\perp^W| = 1/6$.
\end{enumerate}

The result \eqref{eq:k3_w_unification} is the substrate-handle that feeds the Layer 4 V--A coupling derivation in §\ref{sec:sector_a} of the present paper.

\subsection{THEO-SD-CHIR-2 inheritance: qDP/eDP sector closure}
\label{sec:theo_sd_chir_2}

THEO-SD-CHIR-2 \cite{ResearchFrontier, Capotauro} extends the cross-sector unification at substrate level to the qDP/eDP sector (electromagnetic-handedness manifestation under the OPEN-SD-CHIR-PRIMITIVE umbrella). At full Layer 3 rigor:
\begin{equation}
\left| M^{qDP} \right| = \frac{\chi}{6} = \frac{\varphi^{-3}}{6} \approx 0.0394
\label{eq:qDP_chirality}
\end{equation}
identical to $|M^{K_3}|$ and $|M^W|$, completing three-way cross-sector unification at substrate level.

The qDP/eDP sector substrate object is a Linear-ZBW configuration on a $\pm$qCP center at $v_{\text{host}}$. Its stabilizer subgroup is $D_{5d}$ (order 20), obtained via antipodal-pair refinement of the $C_{5v}$ stabilizer at a single first-shell vertex. Its pairing-convention generator $\zeta^{qDP}$ is identified as the combined $CP$ operation \cite{Capotauro, SM2v1} --- host-CP-centered spatial inversion combined with $\hat{n}$-flip combined with qCP-sign flip; the three flips together preserve the substrate's overall chirality content. This is the analog of the W-bracelet's $\zeta^W$ icosahedral-center inversion with an additional charge-conjugation factor that reflects the Capotauro mechanism's three-way coupling (vertex direction $\times$ $\hat{n}$ direction $\times$ qCP-sign) per SM-2 v1.0 §10 \cite{SM2v1}.

The matter-doublet basis under $\sigma_1^{qDP} \zeta^{qDP}$-EVEN pairing convention spans a 2D subspace of $A_{1g} \oplus A_{2u}$ in $D_{5d}$, parallel to the K3 and W-bracelet matter-doublet structures. The chirality operator $\Cchi^{qDP} \in A_{2u}(D_{5d})$ sourced from $\hat{n}$'s pseudoscalar structure has non-vanishing matrix element via Wigner-Eckart, factorizing as amplitude factor $|M_{\text{amp}}^{qDP}| = \chi$ (chirality-eigenvalue matching analog of K3 and W-bracelet) times cage-shell factor $|M_\perp^{qDP}| = \dGamma/\Vcage = 2/12 = 1/6$ via cage-shell averaging on the same icosahedral cage shared across all three sectors.

The result \eqref{eq:qDP_chirality} is the substrate-handle that feeds the Layer 4 chiral-polarity-bias derivation in §\ref{sec:sector_b} of the present paper.

\subsection{The three-way cross-sector unification at substrate level}
\label{sec:three_way_unification}

Combining THEO-CAP-1 \cite{ResearchFrontier, Capotauro} (K3-doublet sector; theorem \#62), THEO-SD-CHIR-1 (W-bracelet sector; theorem \#63), and THEO-SD-CHIR-2 (qDP/eDP sector; theorem \#64) gives the three-way cross-sector unification at substrate level under the OPEN-SD-CHIR-PRIMITIVE umbrella:
\begin{equation}
\boxed{\left| M^{K_3} \right| = \left| M^W \right| = \left| M^{qDP} \right| = \frac{\chi}{6} = \frac{\varphi^{-3}}{6} \approx 0.0394}
\label{eq:umbrella_three_way}
\end{equation}
The structural identity claim at substrate level: three structurally distinct substrate objects (K3-base triangular face, W-bracelet Petrie hexagon, Linear-ZBW configuration) with three structurally distinct stabilizer subgroups ($D_{3d}$ at order 12, $D_6$ at order 12, $D_{5d}$ at order 20) and three structurally distinct pairing-convention generators ($\zeta^{K_3}$ = 3D-inversion through K3-base centroid (purely geometric); $\zeta^W$ = icosahedral-center inversion in 4D ambient (purely geometric); $\zeta^{qDP}$ = combined $CP$ (geometric inversion + $\hat{n}$-flip + qCP-sign flip)) all produce \emph{identical} chirality matrix element magnitude $\chi/6 \approx 0.0394$ via the SAME 12-vertex icosahedral cage at $v_{\text{host}}$ with the SAME effective matter-doublet dimension $\dGamma = 2$. The unification is at substrate level (Layer 3) and is established as theorem-level closure in Capotauro v2.0 \cite{Capotauro}.

The unification is achieved via the three-step closure pattern (Step 1 substrate-locality + Step 2 cage-shell averaging + Step 3 sector-specific pairing convention). This pattern is programme-level technique established by THEO-SD-CHIR-1 and confirmed structurally complete across three manifestations of the OPEN-SD-CHIR-PRIMITIVE umbrella by THEO-SD-CHIR-2 \cite{Capotauro}. The pattern templates the remaining two manifestations (iv) thermodynamic causal arrow and (v) cosmological-vacuum asymmetry under THEO-SD-CHIR-N convention for future work \cite{ResearchFrontier}.

\subsection{What Layer 3 delivered; what Layer 4 must deliver}
\label{sec:layer3_layer4}

Layer 3 (substrate-level closure via Capotauro v2.0) delivers the substrate-handle: identical chirality matrix element magnitude $\chi/6$ on three sectors at substrate cutoff $\Lambdasub = \elledge^{-1}$, with zero free parameters (the magnitude $\chi = \varphi^{-3}$ is derived from substrate primitive 4D direction $\hat{n}$ + 600-cell polytope edge-length ratios via perturbative-distance-ratio constraint; the cage-shell factor $1/6 = \dGamma/\Vcage = 2/12$ is the ratio of two integer-valued topological invariants).

Layer 4 (continuum-EFT projection closure via the present joint paper) must deliver the substrate-handle-to-observable propagation: the magnitude $\chi/6$ at substrate cutoff propagates to continuum-limit effective coupling magnitude at observable kinematic scales $\muobs \ll \Lambdasub$, and then propagates to observable empirical phenomena via standard SM machinery (Yang-Mills EFT for V--A coupling kinematic projection; effective free-energy partition-function framework for chiral-polarity-bias thermodynamic projection). The propagation must be sector-agnostic at the bridge step (§\ref{sec:bridge}) and sector-specific at the kinematic / thermodynamic step (§\ref{sec:sector_a}, §\ref{sec:sector_b}).

The Layer 3 / Layer 4 dichotomy is preserved as architectural framing throughout the paper: Layer 3 results are inherited at theorem-level rigor from Capotauro v2.0 without modification; Layer 4 results are derived in this paper at theorem-level rigor under the conditional-closure framework with foundational input inheritance.

\subsection{Capotauro v2.0 axiom set load-bearing for the joint paper}
\label{sec:axiom_inheritance}

The joint paper's conditional closures rest on the CPP axiom set \cite{AxiomDocument} with most load-bearing axioms inherited from Capotauro v2.0:

\begin{itemize}[leftmargin=*]
\item \textbf{AXIM-1} (CP existence): load-bearing for substrate primitives + FI-CHIR-CONT-1 substrate primitive 4D direction + FI-CHIR-CONT-12 continuum-EFT chirality-projection structure on continuum fermion fields.
\item \textbf{AXIM-2} (600-cell topology): load-bearing for polytope-geometric invariants (edge-length ratios, vertex counts, icosahedral cage structure) + FI-CHIR-CONT-10 + FI-CHIR-CONT-13 sector specializations.
\item \textbf{AXIM-3} (Dipole Sea / DI-bit propagation): load-bearing for FI-CHIR-CONT-11 Yang-Mills EFT framework + FI-CHIR-CONT-14 SM-2 effective free-energy / partition-function framework via dipole-sea propagation mechanism producing continuum gauge structure and thermodynamic structure.
\item \textbf{AXIM-4} (SSV interaction / Nexus): load-bearing for FI-CHIR-CONT-9 Substrate-Locality inheritance and Yang-Mills EFT gauge interaction at continuum level.
\item \textbf{AXIM-7} (Substrate-stress): load-bearing for FI-CHIR-CONT-1 substrate primitive direction sourcing chirality operator under Picture B substrate-orientation-field framework.
\end{itemize}

The foundational input stack inherited from Capotauro v2.0 is FI-C-1 through FI-C-10 \cite{Capotauro} plus the two Reading C foundational inputs FI-C-RC-1 (primitive 4D direction $\hat{n}$) and FI-C-RC-2 (vertex-aligned reading $\hat{n} = v_{\text{host}}$). These are re-labelled FI-CHIR-CONT-1 through FI-CHIR-CONT-9 in the joint paper's foundational input inventory (§\ref{sec:bridge}). Six additional sector-specific foundational inputs (FI-CHIR-CONT-10 through FI-CHIR-CONT-15) are introduced at §\ref{sec:sector_a} and §\ref{sec:sector_b} for the Yang-Mills EFT framework and effective free-energy partition-function framework respectively.

% ============================================================
% §3 BRIDGE — SUBSTANTIVE CONTENT (DRAFTED FROM SKETCH)
% ============================================================
\section{The shared substrate-handle-to-effective-coupling bridge}
\label{sec:bridge}

\begin{remark}[v0.5 SKETCH SCAFFOLDING NOTE]
\textit{This section drafts the §3 bridge content from the comprehensive working sketch at \texttt{flagship\_papers/chirality\_continuum/sketches/substrate\_to\_continuum\_bridge.md} (\textasciitilde{}1156 lines canonical proof chain across Patches 0485+0486+0487). The substantive content is closed at theorem-level rigor as THEO-CHIR-CONT-1 (theorem \#65; Patch 0487; sub-claim (a) closure of joint Layer 4 paper). v0.5 paper polish will integrate sketch content into this section via subsequent patches.}
\end{remark}

[v0.5 placeholder: Section 3 substantive content to be drafted from working sketch \texttt{substrate\_to\_continuum\_bridge.md} during v0.5 paper polish at Patches 0498+. Includes Theorem 3.1 (programme-level registration THEO-CHIR-CONT-1; Substrate-Handle-to-Effective-Coupling Bridge Theorem) with sub-statements Lemma 3.1 = THEO-CHIR-CONT-1.1 (Symmetry-Content Preservation under Continuum-Limit Projection $\Phi$), Theorem 3.2 = THEO-CHIR-CONT-1.2 (Continuum Operator Identification at Sector-Agnostic Level), and Theorem 3.3 = THEO-CHIR-CONT-1.3 (Magnitude Inheritance via Topological Projection). Definition 11.2.1 (continuum-limit projection map $\Phi$ via Wilson-Fisher block-spin renormalization), Definition 15.1.1 (Topological Substrate Quantity), Claims 15.1.2 (topological-substrate-quantity character of $|\chi|$) and 15.2.1 (topological character of cage-shell factor $1/6 = \dGamma/\Vcage$).]

% ============================================================
% §4 SECTOR A — SUBSTANTIVE CONTENT (DRAFTED FROM SKETCH)
% ============================================================
\section{Sector A: Electroweak V--A Coupling Derivation}
\label{sec:sector_a}

\begin{remark}[v0.5 SKETCH SCAFFOLDING NOTE]
\textit{This section drafts the §4 Sector A content from the comprehensive working sketch at \texttt{flagship\_papers/chirality\_continuum/sketches/sector\_a\_va\_coupling.md} (\textasciitilde{}885 lines canonical proof chain across Patches 0488+0489+0490+0491). The substantive content is closed at theorem-level rigor as THEO-CHIR-CONT-2 (theorem \#66; Patch 0491; §B sub-claim (b)+(c)+(d)+(e) closure). v0.5 paper polish will integrate sketch content into this section via subsequent patches.}
\end{remark}

[v0.5 placeholder: Section 4 substantive content to be drafted from working sketch \texttt{sector\_a\_va\_coupling.md} during v0.5 paper polish at Patches 0498+. Includes Theorem 4.4 (programme-level registration THEO-CHIR-CONT-2; Sector A Yang-Mills EFT V--A Coupling Derivation Theorem) with four sub-claim consequences: (b) sector-specific operator identification via three structural identifications ($\zeta^{\text{cont,W}} \leftrightarrow \gamma_5$, matter-doublet $\leftrightarrow \{\psi_R, \psi_L\}$, operator $\leftrightarrow$ V--A current); coupling pinned to $g^V_{LL} = 1$. (c) Michel parameter $\rho = 3/4$ at finite mass via textbook V--A four-fermion kinematics; PDG 2024 $\rho^{\text{obs}} = 0.7497 \pm 0.0010$ within $0.3\sigma$. (d) 100\% LH at massless helicity limit via chirality-helicity coincidence; multi-sector empirical validation (Goldhaber 1958, LEP/SLC $\tau$-polarization, LHC top-quark spin-correlation, modern neutrino constraints). (e) Capotauro Falsifier 6 ACTIVATION with three thresholds (A) Michel + (B) massless-helicity + (C) leptogenesis CP-asymmetry --- sharpest direct test of substrate-handle magnitude inheritance bypassing kinematic intermediaries.]

% ============================================================
% §5 SECTOR B — SUBSTANTIVE CONTENT (DRAFTED FROM SKETCH)
% ============================================================
\section{Sector B: Chiral-Polarity-Bias Derivation}
\label{sec:sector_b}

\begin{remark}[v0.5 SKETCH SCAFFOLDING NOTE]
\textit{This section drafts the §5 Sector B content from the comprehensive working sketch at \texttt{flagship\_papers/chirality\_continuum/sketches/sector\_b\_chiral\_polarity\_bias.md} (\textasciitilde{}967 lines canonical proof chain across Patches 0492+0493+0494+0495). The substantive content is closed at theorem-level rigor as THEO-CHIR-CONT-3 (theorem \#67; Patch 0495; §C sub-claim (f)+(g)+(h)+(i) closure). v0.5 paper polish will integrate sketch content into this section via subsequent patches.}
\end{remark}

[v0.5 placeholder: Section 5 substantive content to be drafted from working sketch \texttt{sector\_b\_chiral\_polarity\_bias.md} during v0.5 paper polish at Patches 0498+. Includes Theorem 5.5 (programme-level registration THEO-CHIR-CONT-3; Sector B SM-2 Effective Free-Energy Chiral-Polarity-Bias Exclusion Theorem) with four sub-claim consequences: (f) sector-specific operator identification $\OeffqDP \leftrightarrow \DeltaFqDP$ via three structural identifications ($\zeta^{\text{cont,qDP}} \leftrightarrow$ combined $CP$, matter-doublet $\leftrightarrow \{|\text{LZBW},+\rangle, |\text{LZBW},-\rangle\}$, operator $\leftrightarrow \DeltaFqDP = F[\text{LZBW},+] - F[\text{LZBW},-]$). (g) Substrate-level stabilization energy $|M^{\text{eff,qDP}}| = \chi/6 \approx 0.0394$ via THEO-CHIR-CONT-1.3 topological-projection + THEO-SD-CHIR-2 Finding C-W46 inheritance + sector-agnostic claim verification + sub-leading correction quantification. (h) Exclusion bound at observable thermodynamic scales $\DeltapLR \approx \chi/6 \approx 0.0394$ via Boltzmann-like thermodynamic distribution + PRED-O-25 inheritance + BAU back-derivation empirical anchor (Davidson, Nardi, Nir 2008). (i) SM cross-validation via three-track cross-validation against SM-2 v1.0 §10 + §4 THEO-CHIR-CONT-2 + joint paper cross-sector unification framing at observable scale.]

% ============================================================
% §6 CROSS-SECTOR UNIFICATION — SUBSTANTIVE CONTENT (DRAFTED FROM SKETCH)
% ============================================================
\section{Cross-Sector Unification: The Structural Identity Claim}
\label{sec:cross_sector_unification}

\begin{remark}[v0.5 SKETCH SCAFFOLDING NOTE]
\textit{This section drafts the §6 cross-sector unification synthesis content from the working sketch at \texttt{flagship\_papers/chirality\_continuum/sketches/cross\_sector\_unification.md} (\textasciitilde{}500+ lines across Patch 0496). The substantive content is synthesis of themes that emerged through §3 + §4 + §5 closures (THEO-CHIR-CONT-1+2+3). v0.5 paper polish will integrate sketch content into this section via subsequent patches.}
\end{remark}

[v0.5 placeholder: Section 6 substantive content to be drafted from working sketch \texttt{cross\_sector\_unification.md} during v0.5 paper polish at Patches 0498+. Six sub-sections: (6.1) Shared substrate handle: $|M^{K_3}| = |M^W| = |M^{qDP}| = \chi/6 \approx 0.0394$ across three sectors via SAME 12-vertex icosahedral cage + SAME matter-doublet dimension 2. (6.2) OPEN-SD-CHIR-PRIMITIVE umbrella: 3 of 5 manifestations at full Layer 4 rigor under THEO-CHIR-CONT-N convention; manifestations (iv)+(v) at substrate-level only pending future work. (6.3) Second cross-sector closure pattern in CPP after SF-4 v4.0; methodological pattern strengthening. (6.4) Structural identity claim of the paper: single substrate primitive $\hat{n}$ + $|\chi| = \varphi^{-3}$ controls every parity-sensitive observable via shared three-step machinery; zero free parameters. (6.5 NEW) Joint paper format structural efficiency validation: combined §3+§4+§5 closure cost 11 sessions vs Venue (b) fallback 15--22 sessions; savings 4--11 sessions; cross-sector convergence at observable level as structural prediction rather than emergent empirical coincidence. (6.6 NEW) v0.5 SHIP readiness assessment: READY for paper polish + reviewer cycle.]

% ============================================================
% §7 PREDICTIONS AND FALSIFIERS — SCAFFOLD
% ============================================================
\section{Predictions and Falsifiers}
\label{sec:predictions}

\begin{remark}[v0.5 SCAFFOLD NOTE]
\textit{This section will be substantively drafted at v0.5 paper polish per v0.1 outline §1.9 specification \cite{ResearchFrontier}.}
\end{remark}

[v0.5 placeholder: Section 7 substantive content to be drafted at Patches 0498+ paper polish from v0.1 outline §1.9. Includes Table of zero-parameter predictions (substrate-level chirality magnitude $\chi/6 \approx 0.0394$ inherited from Capotauro v2.0; V--A coupling strength 75\% LH at finite mass / 100\% LH at massless limit; Michel parameter $\rho = 3/4$ matching PDG 2024 $\rho^{\text{obs}} = 0.74979 \pm 0.00026$ at $10^{-4}$ level; Linear-ZBW-on-$-$qCP stabilization energy producing $\DeltapLR \approx 0.0394$ matching BAU back-derivation $\sim 0.04$ at 2\% level; cross-sector unification identity at substrate level inherited from Capotauro v2.0). Six falsifiers per v0.1 outline §1.9 §7.2: (1) Layer 4 EFT projection failure (Capotauro Falsifier 6 activation per §4); (2) Michel parameter deviation $\rho \neq 3/4$ beyond PDG precision; (3) right-handed neutrino observation coupling to $W^\pm$; (4) positive down-type quark observation at SM-accessible scales; (5) cosmological constraint on chiral-polarity-bias from BBN/CMB; (6) cross-sector unification breakdown via differential scaling with substrate magnitude.]

% ============================================================
% §8 OPEN THEOREM-LEVEL WORK — SCAFFOLD
% ============================================================
\section{Open Theorem-Level Work}
\label{sec:open_work}

\begin{remark}[v0.5 SCAFFOLD NOTE]
\textit{This section will be substantively drafted at v0.5 paper polish per v0.1 outline §1.10 specification.}
\end{remark}

[v0.5 placeholder: Section 8 substantive content to be drafted at Patches 0498+ paper polish. Includes: (8.1) Sub-claims under sectoral closures (post-v1.0 SHIP work) --- Layer 1 substrate-dynamics derivation of $\hat{n}$ tied to Q1$'$+Q1$'$.A Layer 3 promotion; SM-2 v2.0+ chiral-polarity-bias mechanism from CPP axioms A3 + A7. (8.2) OPEN-SD-CHIR-PRIMITIVE umbrella future-window work --- manifestation (iv) thermodynamic causal arrow + (v) cosmological-vacuum asymmetry via THEO-CHIR-CONT-4/-5 candidates templated by THEO-CHIR-CONT-N convention. (8.3) Cross-validation candidates --- Route (ii) substrate-mechanism via Patch 0367 $W^0$ neutrino scattering centroid-decoupling sketch; future-window work pending Layers A/B/C operational development. (8.4) Picture A alternative continuum-EFT framework parametrization (OPEN-FP-SF-4-1 candidate); complementary closure to Picture B Wigner-Eckart EFT framework completed in this paper.]

% ============================================================
% §9 DISCUSSION — SCAFFOLD
% ============================================================
\section{Discussion}
\label{sec:discussion}

\begin{remark}[v0.5 SCAFFOLD NOTE]
\textit{This section will be substantively drafted at v0.5 paper polish per v0.1 outline §1.11 specification.}
\end{remark}

[v0.5 placeholder: Section 9 substantive content to be drafted at Patches 0498+ paper polish. Discusses: (9.1) Joint paper format as established CPP methodology --- second cross-sector closure after SF-4 v4.0; templates future cross-sector closure work under OPEN-SD-CHIR-PRIMITIVE umbrella and other umbrellas. (9.2) THEO-CHIR-CONT-N convention as programme-level template for Layer 4 continuum-EFT projection closures under any cross-sector unification umbrella. (9.3) Structural identity claim of the paper as confirmation of substrate-primitive parsimony in CPP framework --- one 4D direction $\hat{n}$ + one polytope-geometric magnitude $\varphi^{-3}$ controls multiple seemingly-disparate Standard Model phenomenological features. (9.4) Forward trajectory: future-window Layer 4 closures under manifestations (iv)+(v); Picture A alternative complementary closure; FI-CHIR-CONT-1/2 Layer 3 first-principles promotion (Q1$'$+Q1$'$.A). (9.5) Implications for the broader CPP programme architecture: cross-sector closure as a recurring closure pattern at scale within the conditional-closure-framework.]

% ============================================================
% §10 REFERENCES
% ============================================================
\begin{thebibliography}{99}

\bibitem{Capotauro}
T. L. Abshier, ``The Capotauro Mechanism: Chirality on the K3-Doublet from Substrate-Vacuum Broken-Symmetry Physics,'' Conscious Point Physics Flagship Paper Series, Version 2.0 v1.0 SHIPPED, 19 May 2026. \texttt{flagship\_papers/capotauro/capotauro.tex}.

\bibitem{AbshierCPP2026}
T. L. Abshier, \emph{Conscious Point Physics: A Theoretical Framework}, Hyperphysics Institute, 2026. Master programme reference document.

\bibitem{ResearchFrontier}
Conscious Point Physics Programme, ``Research Frontier and Theorem Registry,'' continuously updated. \texttt{research\_frontier.md} + \texttt{theorem-registry.md} (Hyperphysics-Institute/CPP GitHub repository).

\bibitem{SF2v1}
T. L. Abshier, ``Standard Model Sector Foundational Paper SF-2: Electroweak Sector via 600-Cell Substrate Symmetry,'' Conscious Point Physics Flagship Paper Series, v1.0 SHIPPED, 14 May 2026. \texttt{flagship\_papers/electroweak/sf-2\_electroweak.tex}.

\bibitem{SM2v1}
T. L. Abshier, ``Standard Model Paper SM-2: Mass Generation and Geometric Hierarchies in Conscious Point Physics,'' Conscious Point Physics SM Series, v1.0, 2026. \texttt{series\_standard\_model/papers/SM-2\_mass\_generation\_geometric\_hierarchies.tex}.

\bibitem{EW5}
T. L. Abshier, ``EW-5: Yang-Mills Continuum-EFT Theorem,'' Conscious Point Physics EW Series, 2026. THEO-EW-8 thm:YM\_EFT proof outline inherited at FI-CHIR-CONT-11.

\bibitem{ConditionalClosure}
Conscious Point Physics Programme, ``Conditional Closure Framework for Theorem Registry,'' templates/conditional\_closure\_framework.md. Programme-level methodology document.

\bibitem{AxiomDocument}
T. L. Abshier, \emph{The Axioms of Conscious Point Physics: AXIM-1 through AXIM-7}, Hyperphysics Institute, 2026. Programme-level foundational document.

\bibitem{ReviewerNotes}
Reviewer letters for Capotauro v2.0 v0.9 / v0.9.1 cycle. \texttt{flagship\_papers/capotauro/reviews/} (ChatGPT round-2 + CoPilot v0.9 + Grok v0.9; Session 135 Patches 0474--0476).

\bibitem{DavidsonNardiNir}
S. Davidson, E. Nardi, and Y. Nir, ``Leptogenesis,'' \emph{Physics Reports} \textbf{466}, 105 (2008). arXiv:0802.2962 [hep-ph].

\bibitem{Wu1957}
C. S. Wu, E. Ambler, R. W. Hayward, D. D. Hoppes, and R. P. Hudson, ``Experimental Test of Parity Conservation in Beta Decay,'' \emph{Physical Review} \textbf{105}, 1413 (1957).

\bibitem{Goldhaber1958}
M. Goldhaber, L. Grodzins, and A. W. Sunyar, ``Helicity of Neutrinos,'' \emph{Physical Review} \textbf{109}, 1015 (1958).

\bibitem{CommonsBucksbaum}
E. D. Commins and P. H. Bucksbaum, \emph{Weak Interactions of Leptons and Quarks}, Cambridge University Press, 1983.

\bibitem{ChengLi}
T.-P. Cheng and L.-F. Li, \emph{Gauge Theory of Elementary Particle Physics}, Oxford University Press, 1984.

\bibitem{Peskin}
M. E. Peskin and D. V. Schroeder, \emph{An Introduction to Quantum Field Theory}, Westview Press, 1995. Especially §3.4 (chirality and helicity for fermion fields) and §17.2 (electroweak charged current).

\bibitem{TWISTrho}
TWIST Collaboration, ``Measurement of the muon decay parameter $\rho$,'' \emph{Physical Review Letters} \textbf{106}, 041804 (2011); PDG 2024 muon decay parameters.

\bibitem{Michel1950}
L. Michel, ``Interaction between four half-spin particles and the decay of the $\mu$-meson,'' \emph{Proceedings of the Physical Society A} \textbf{63}, 514 (1950).

\bibitem{MarcianoSirlin1988}
W. J. Marciano and A. Sirlin, ``Radiative corrections to neutrino-induced neutral current phenomena in the $SU(2)_L \times U(1)$ theory,'' \emph{Physical Review D} \textbf{22}, 2695 (1980); ``Electroweak radiative corrections to $\tau$ decay,'' \emph{Physical Review Letters} \textbf{61}, 1815 (1988).

\bibitem{LEPElectroweak}
ALEPH, DELPHI, L3, OPAL, SLD Collaborations, ``Precision electroweak measurements on the $Z$ resonance,'' \emph{Physics Reports} \textbf{427}, 257 (2006). $\tau$-polarization measurements.

\bibitem{ATLASTopSpin}
ATLAS Collaboration, ``Measurements of top-quark pair spin correlation observables at $\sqrt{s} = 13$ TeV,'' \emph{European Physical Journal C} \textbf{80}, 754 (2020).

\bibitem{CMSTopSpin}
CMS Collaboration, ``Measurement of the top-quark polarization at $\sqrt{s} = 13$ TeV,'' \emph{Physical Review D} \textbf{100}, 072002 (2019).

% Additional references to be added at v0.5 paper polish:
% - PDG 2024 (Review of Particle Physics, Workman et al. 2024)
% - SF-4 v4.0 SHIP paper (THEO-SF-4-5 first-cross-sector-closure precedent)
% - Adler-Bardeen anomaly nonrenormalization theorem
% - Wilson-Fisher block-spin renormalization references
% - Atiyah-Singer index theorem references
% - Polytope geometry / 600-cell references (Coxeter)

\end{thebibliography}

% ============================================================
% END OF DOCUMENT
% ============================================================
\end{document}
